\documentclass[11pt,a4paper]{article}

\usepackage[margin=1in]{geometry}
\usepackage{amsmath, amssymb}
\usepackage{graphicx}
\usepackage{hyperref}
\usepackage{booktabs}
\usepackage{siunitx}
\usepackage{enumitem}
\usepackage{xcolor}
\usepackage[normalem]{ulem}

\hypersetup{colorlinks=true, linkcolor=blue!50!black, urlcolor=blue!50!black, citecolor=blue!50!black}

\title{Basal mixing model: forward dynamics and sampling strategies}
\author{BasalMixing repository}
\date{\today}

\newcommand{\Kr}{\ensuremath{{}^{81}\mathrm{Kr}}}
\newcommand{\Ar}{\ensuremath{\delta^{40}\mathrm{Ar}_{\mathrm{atm}}}}
\newcommand{\tpred}{\ensuremath{t_{\mathrm{pred}}}}
\newcommand{\teq}{\ensuremath{t_{\mathrm{eq}}}}
\newcommand{\code}[1]{\texttt{#1}}

\begin{document}

\maketitle

\begin{abstract}
\noindent
This document describes the 1-D basal mixing model used in the BasalMixing
repository to interpret deep-ice \Kr{} and \Ar{} measurements at GISP2, and the
three Markov-chain Monte Carlo (MCMC) strategies the code provides for
inferring the mixing parameters. The three strategies (\code{:profile},
\code{:sampled}, \code{:marginal}) differ in how they handle a single nuisance
quantity: the time the mixing system has been running before the measurement.
We argue that the marginal-likelihood strategy is the correct Bayesian
treatment of this problem and is what should be used for publication-quality
inference; the other two strategies are kept for backward compatibility and
diagnostic comparison.
\end{abstract}

\tableofcontents

\section{The basal mixing model}
\label{sec:model}

\subsection{Physical setting}

GISP2 retrieves ice from the bottom of the Greenland Ice Sheet. The deepest
$\sim$13~m of the core sits above the bedrock and, unlike the rest of the
column, is thought to have been mechanically and diffusively mixed with
material below. Specifically: above $z=z_{\mathrm{lim}}\approx 3040\,\mathrm{m}$
the ice is treated as undisturbed (``clean'' ice); below $z_{\mathrm{lim}}$,
ice is increasingly contaminated by debris-rich (``dirty'') material that
acts as a sink/source for the tracers we measure. The model treats this lower
layer as a 1-D advection--diffusion column extending to a fixed bedrock at
$z_{\mathrm{bed}}=3053.44\,\mathrm{m}$.

\subsection{State variables and grid}

The model carries two state variables at every cell centre $z_j$ in a finite-volume
discretization:

\begin{itemize}
  \item $c_{81}(z, t)$: relative \Kr{} concentration, dimensionless. Initialized
        to 1 at $t=0$ (modern abundance) at every depth.
  \item $c_{40}(z, t)$: \Ar{} concentration, in $\mathrm{cc}\,\mathrm{m}^{-3}$.
        Initialized to the modern equilibrium value
        $A_{40}^{0} = (100\,\mathrm{cm})^{3}\cdot a_{\mathrm{tac}}\cdot
         f_{\mathrm{atm}} \approx 8000\,\mathrm{cc/m^{3}}$,
        where $a_{\mathrm{tac}}=0.08$ is the typical air content and
        $f_{\mathrm{atm}}=0.00934$ is the present-day atmospheric \Ar{} volume
        fraction.
\end{itemize}

The model integrates from $t=0$ (initial condition) up to some upper time
$t_1$ (typically $3000\,\mathrm{kyr}$), advancing with a forward Euler step
$dt=0.2\,\mathrm{kyr}$ (with an automatic retry at $dt/2$ if the integrator
produces non-physical, negative concentrations).

\subsection{Mixing rate profile}

The depth-dependent mixing rate is parameterised by two values
$m_{\mathrm{clean}}$ and $m_{\mathrm{dirty}} = f_{\mathrm{dirty}}\,m_{\mathrm{clean}}$
joined by a smooth transition of thickness $\delta$ centred at
$z_{\mathrm{lim}}$. Using $w_1$ and $w_2$ for the two tanh-shaped step functions,
\begin{align}
w_1(z) &= \tfrac{1}{2}\!\left(1 + \tanh\!\big(\tfrac{\beta}{\delta}(z - z_{\mathrm{lim}})\big)\right), \\
w_2(z) &= \tfrac{1}{2}\!\left(1 + \tanh\!\big(\tfrac{\beta}{\delta}(z - z_{\mathrm{lim}} - \delta)\big)\right), \\
\Phi(z) &= m_{\mathrm{dirty}}\,w_1(z)\,w_2(z) + m_{\mathrm{clean}}\,w_1(z)\,(1 - w_2(z)),
\end{align}
with sharpness $\beta=50$ (essentially a step on the discretisation scale).
So $\Phi=0$ in clean ice, $\approx m_{\mathrm{clean}}$ in the transition zone of
width $\delta$, and $\approx m_{\mathrm{dirty}}$ in the dirty zone below.

The diffusivity is $D(z) = \Phi(z)\,L_{\mathrm{ref}}$ with $L_{\mathrm{ref}}=1\,\mathrm{m}$
treated as a reference length (it just rescales $m_{\mathrm{clean}}$; we keep
$L_{\mathrm{ref}}=1$ and let $m_{\mathrm{clean}}$ carry the magnitude).

\subsection{Governing equations}

The continuum dynamics of $c\in\{c_{81}, c_{40}\}$ are
\begin{equation}
\frac{\partial c}{\partial t} = \frac{\partial}{\partial z}\!\left(D(z)\,\frac{\partial c}{\partial z}\right) + S_c(z, t),
\label{eq:pde}
\end{equation}
discretised cell-centre to cell-centre. The source/sink term $S_c$ differs by
tracer:

\begin{itemize}
  \item For \Kr{}, $S_{81}(z, t) = -\lambda\, c_{81}(z, t)$ with
        $\lambda = \ln 2 / t_{1/2}$ and $t_{1/2}=229\,\mathrm{kyr}$, except in
        the clean-ice cells for times $t > t_{\mathrm{old}}$, where the decay
        term is switched off. (Physically: the clean ice was deposited at most
        $t_{\mathrm{old}}$ ago and so its \Kr{} has decayed for at most that
        long.)
  \item For \Ar{}, $S_{40} = 0$ in the interior; at the bedrock cell only,
        a flux $F_{40}$ is added: $S_{40}(z_{\mathrm{bed}}, t) =
        F_{40}/\Delta z_{\mathrm{bed}}$.
\end{itemize}

The observable quantities derived from the state are
\begin{align}
\tau_{81}(z, t) &= -\,t_{1/2}\,\log_2\!\big(c_{81}(z, t)\big),
 \qquad \text{Kr-81 closed-system age, kyr,} \\
\delta^{40}\mathrm{Ar}_{\mathrm{atm}}(z, t) &= 1000\,\left(\frac{c_{40}(z,t)}{A_{40}^{0}} - 1\right) - 0.066\,t,
 \qquad \text{$\delta^{40}\mathrm{Ar}$, \textperthousand},
\end{align}
where the $-0.066\,t$ term in the latter expression corrects for atmospheric
\Ar{} growth over time (Bender et al. 2010).

\subsection{Parameters}

The five sampled parameters and the two fixed observation-noise scales are:

\begin{center}
\small
\begin{tabular}{lllp{4.5cm}l}
\toprule
Symbol & Meaning & Units & Prior & Code name \\
\midrule
$\delta$ & transition-zone thickness & m & $\mathrm{Uniform}(0.3,\,2.0)$ & \code{delta} \\
$m_{\mathrm{clean}}$ & mixing rate in transition zone & m/kyr & $\mathrm{Normal}(0.03,\,0.002)$ trunc.\ $\geq 0$ & \code{m\_clean} \\
$f_{\mathrm{dirty}}$ & dirty-to-clean mixing-rate ratio & --- & $\mathrm{Uniform}(4.0,\,6.5)$ & \code{f\_dirty} \\
$t_{\mathrm{old}}$ & clean-ice maximum age & kyr & $\mathrm{Normal}(250,\,25)$ trunc.\ $\geq 0$ & \code{t\_old} \\
$F_{40}$ & bedrock \Ar{} source flux & cc/m$^2$/kyr & $\mathrm{Uniform}(0.004,\,0.007)$ & \code{F\_ar40} \\
\midrule
$\sigma_{81}$ & \Kr{} age obs.\ noise (\emph{fixed}) & kyr & 30 & \code{$\sigma$\_k81} \\
$\sigma_{40}$ & $\delta^{40}\mathrm{Ar}$ obs.\ noise (\emph{fixed}) & \textperthousand & 0.03 & \code{$\sigma$\_dar40} \\
\bottomrule
\end{tabular}
\end{center}

The choice to fix $\sigma_{81}$ and $\sigma_{40}$ rather than place priors on
them was made to improve sampling efficiency. Both values are close to the
analytical errors of the measurements; relaxing them to half-Cauchy or
exponential priors is a minor change to the model if needed.

\subsection{Observations}

The model is compared against two data sets, both at observation depths
$z_i$ corresponding to specific GISP2 ice samples:

\begin{itemize}
  \item Three \Kr{} closed-system ages, $\tau_{81}^{\mathrm{obs}}$, at depths
        $3045$--$3050\,\mathrm{m}$ (typically reported as ages of $577$, $649$, and
        $714\,\mathrm{kyr}$).
  \item Ten \Ar{} measurements from Bender et al.\ (2010), spanning depths
        $3036.5$ to $3052.4\,\mathrm{m}$, with values rising from
        $\sim 0$ \textperthousand{} in the clean zone to
        $\sim 0.5$ \textperthousand{} near bedrock.
\end{itemize}

For depths between cell centres the model predictions are linearly
interpolated. Both data sets are assumed to be independent Gaussian
observations of the corresponding model quantities at the depths in question.

\section{The likelihood and the time-of-comparison problem}
\label{sec:tpred}

For given parameters $\theta = (\delta, m_{\mathrm{clean}}, f_{\mathrm{dirty}},
t_{\mathrm{old}}, F_{40})$ and integration time $t$, the model produces
predictions $\hat{\tau}_{81}(z_i;\theta,t)$ and $\widehat{\delta^{40}\mathrm{Ar}}(z_j;\theta,t)$
at the observation depths $z_i,z_j$. Conditional on $\theta$ and $t$, the
log-likelihood is the usual Gaussian:
\begin{equation}
\log L(t \mid \theta, \mathbf{y}) = -\tfrac{1}{2\sigma_{81}^2}\sum_{i}\!\left(\hat{\tau}_{81}(z_i;\theta,t) - y_i^{\,81}\right)^{2}
- \tfrac{1}{2\sigma_{40}^2}\sum_{j}\!\left(\widehat{\delta^{40}\mathrm{Ar}}(z_j;\theta,t) - y_j^{\,40}\right)^{2} + C,
\label{eq:loglik}
\end{equation}
where $C$ collects $\theta$- and $t$-independent normalising constants.

The wrinkle is that we have no direct information about the integration time
$t$ that should be used to compare the model with the measurements. Physically,
$t$ is the amount of time the mixing process at the base of the ice sheet has
been operating before the present day. The data alone don't pin this number
down: once the mixing column has run long enough to reach a steady state, all
later times look identical, so any $t \geq \teq(\theta)$ produces the same
predictions and therefore the same likelihood.

This non-identifiability is the central reason the three sampling strategies
below differ. They are all valid pieces of inference machinery but answer
different questions.

\section{Sampling strategies}
\label{sec:strategies}

The repository implements three strategies, selectable by the
\code{model\_kind} variable at the top of
\code{run\_basalmixing\_ensemble.jl}. All three use the same affine-invariant
ensemble sampler (\code{Emcee}; Goodman \& Weare 2010) with 32 walkers, so
their wall-clock costs are directly comparable per integration of the forward
model.

\subsection{\code{:profile} --- profile likelihood}

\paragraph{What it does.}
For every parameter proposal $\theta$, integrate the forward model from $t=0$
to $t_1=3000\,\mathrm{kyr}$ on a 1-kyr grid. At each grid point compute the
log-likelihood (\ref{eq:loglik}). Select
\begin{equation}
\tpred^{\star}(\theta) = \arg\max_{t}\log L(t \mid \theta, \mathbf{y}),
\end{equation}
and use $\log L^{\star}(\theta) = \log L(\tpred^{\star} \mid \theta, \mathbf{y})$
as the contribution to the MCMC log-posterior.

\paragraph{Assumptions.}
\begin{itemize}
  \item The ``correct'' time $\tpred$ for each $\theta$ is treated as a
        deterministic function of $\theta$ chosen to make the data fit best.
  \item No prior is placed on $\tpred$.
\end{itemize}

\paragraph{What you get out.}
A posterior on $\theta$ from a non-standard (non-Bayesian) likelihood: the
\emph{profile likelihood} $L^{\star}(\theta)$. The reported $\tpred^{\star}$
per chain sample is the maximizer, not a sample from a posterior.

\paragraph{Pros and cons.}
Computationally efficient because the inner integration can early-terminate
once predicted ages exceed the range of interest. Single-pass per sample,
no extra dimension to mix over. \textit{However:} the posterior on $\theta$
is biased --- it over-credits each $\theta$ for the best $t$ available rather
than averaging the data's support over plausible $t$. The reported $\tpred$
``posterior'' is meaningless as a probability distribution; it's a histogram
of point optima.

\subsection{\code{:sampled} --- explicit \tpred{} prior, joint $(\theta, \tpred)$ posterior}

\paragraph{What it does.}
Treat $\tpred$ as an additional sampled parameter, with a uniform prior
$\tpred \sim \mathrm{Uniform}(700, 3000)\,\mathrm{kyr}$. For each $(\theta,
\tpred)$ proposal, run the forward model once from $t=0$ to $t=\tpred$ and
evaluate (\ref{eq:loglik}) at $t = \tpred$.

\paragraph{Assumptions.}
\begin{itemize}
  \item $\tpred$ is a meaningful physical quantity that we have prior belief
        about (uniform on $[700, 3000]$).
  \item The integration finish time \emph{is} the comparison time.
\end{itemize}

\paragraph{What you get out.}
A proper joint Bayesian posterior on $(\theta, \tpred)$. The marginal on
$\theta$ accounts for uncertainty in $\tpred$ by averaging; the marginal on
$\tpred$ is the data's combined credibility about ``how long has this been
running.''

\paragraph{Pros and cons.}
Bayesianly sound at the level of the joint posterior, and on average
\textit{cheaper} per evaluation than \code{:profile} because most proposed
$\tpred$ are less than 3000 (typical mean $\tpred\approx 1500\,\mathrm{kyr}$,
so a $\sim 2\times$ speed-up).

\textbf{But the posterior on $\tpred$ is essentially unidentifiable.} Once
the system has equilibrated for the proposed $\theta$, every later $t$ gives
the same likelihood; the data don't choose between $\tpred=1500$ and
$\tpred=2500$ when $\teq(\theta)=1000$. The marginal $p(\tpred \mid
\mathbf{y})$ in the chain therefore:
\begin{itemize}
  \item Has a tall peak near the smallest equilibration time that is
        compatible with the data (typically $\sim$1000 kyr), where the
        likelihood first plateaus;
  \item Rises monotonically toward the upper prior bound $t_1=3000\,\mathrm{kyr}$.
        This rise is the mixture-of-uniforms effect: chain samples with
        smaller $\teq(\theta)$ contribute mass to all $t$ above their
        $\teq$, so the higher $t$ values accumulate contributions from
        more chain samples. It is a mathematically correct posterior under
        a $\mathrm{Uniform}(700, 3000)$ prior, but the precise rise rate
        and the location of the upper edge are determined by the
        \emph{arbitrary} choice of $t_1$, not by the data.
\end{itemize}
If one extended the prior to $\mathrm{Uniform}(700, T_{\max})$ for some
$T_{\max} > 3000$, the same rising shape would extend to $T_{\max}$ and pile
up there. The boundary is artificial.

\subsection{\code{:marginal} --- treat \tpred{} as a nuisance parameter, integrate it out}

\paragraph{What it does.}
Recognise that $\tpred$ is not a quantity we actually care about (it has no
independent physical observation) and is unidentifiable past
$\teq(\theta)$. Treat it as a nuisance parameter and marginalise it out of
the likelihood:
\begin{equation}
\log p(\mathbf{y} \mid \theta) \;=\; \log\!\int_{t_0}^{t_1}\!\! L(t \mid \theta, \mathbf{y})\,p(t)\,dt.
\label{eq:marginal-loglik}
\end{equation}
With a uniform prior $p(t) = 1/(t_1 - t_0)$ on $[t_0, t_1]$, the normalisation
is constant in $\theta$ and can be absorbed into the chain's normalising
constant. Numerically we evaluate (\ref{eq:marginal-loglik}) as a trapezoidal
sum on the 1-kyr prediction grid:
\begin{equation}
\log p(\mathbf{y} \mid \theta) \;\approx\; \log\!\left(\sum_{k}\exp\!\big(\log L(t_k \mid \theta)\big)\,\Delta t\right),
\label{eq:marginal-trapz}
\end{equation}
computed via the standard \code{logsumexp} trick to avoid underflow.

\paragraph{Equilibration speed-up.}
For typical $\theta$ the integration spends most of its wall-clock in the
post-equilibration phase, where the integrand is constant. We detect
equilibration online by comparing the predicted observation vectors $W$
slices apart (default $W=100\,\mathrm{kyr}$) and declaring equilibrium when
both \Kr{} and \Ar{} predictions have moved less than tight tolerances
($\sim 1\%$ of the corresponding observation noise scale). When equilibrium
is declared at slice $k_{\mathrm{eq}}$ with time $\teq$ and log-likelihood
$\log L_{\mathrm{eq}}$, the remainder of the integral is added in closed
form:
\begin{equation}
\log p(\mathbf{y} \mid \theta) \;\approx\; \log\!\left(\sum_{k \leq k_{\mathrm{eq}}}\!\exp\!\big(\log L(t_k)\big)\,\Delta t \;+\; (t_1 - \teq)\,\exp\!\big(\log L_{\mathrm{eq}}\big)\right).
\label{eq:marginal-eq}
\end{equation}
In practice the integrator equilibrates around $t \in [1\,000, 1\,500]$ kyr,
giving roughly a 2--3$\times$ speed-up per evaluation.

\paragraph{Assumptions.}
\begin{itemize}
  \item $\tpred$ is a nuisance parameter --- we are not interested in
        sampling its distribution.
  \item The prior on $\tpred$ is uniform over $[t_0, t_1]$. The choice of
        $t_1$ enters the log-likelihood only as an additive constant
        (the normalisation $-\log(t_1-t_0)$), which factors out of the
        chain's relative weights; \emph{the posterior on $\theta$ does not
        depend on $t_1$}, provided $t_1$ is large enough to be past
        equilibration for the bulk of the posterior support. This is the
        key qualitative improvement over \code{:sampled}.
  \item Equilibration is detected within tolerance $\sim 1\%$ of
        $\sigma_{81}$ and $\sigma_{40}$; if for some $\theta$ the system
        equilibrates more slowly than the integration grid runs, the
        explicit sum (\ref{eq:marginal-trapz}) is used and remains correct.
\end{itemize}

\paragraph{What you get out.}
A clean marginal posterior on $\theta$ (5-dimensional, vs.\ 6-dimensional
for \code{:sampled}). \textbf{There is no posterior on $\tpred$ in the chain.}
If a posterior on $\tpred$ is needed for diagnostics, it can be reconstructed
\textit{post hoc} from the chain:
\begin{equation}
p(\tpred \mid \mathbf{y}) \;\propto\; \int p(\tpred \mid \mathbf{y},\theta)\,p(\theta \mid \mathbf{y})\,d\theta
                          \;\approx\; \frac{1}{N}\sum_{s=1}^{N}\frac{L(\tpred \mid \theta^{(s)})}{\int L(t \mid \theta^{(s)})\,dt}\!,
\end{equation}
which averages the per-slice likelihood curves computed during sampling. The
result is meaningful (it summarises what the data say about $\tpred$ in light
of the parameter posterior) without injecting the boundary-pile-up artefact
of \code{:sampled}.

\paragraph{Pros and cons.}
This is the correct Bayesian treatment of an unconstrained nuisance
parameter and is the recommended default. Wall-clock per evaluation is
comparable to \code{:profile} (the integrator runs to $\teq$ rather than to
$t_1$); the chain is one dimension smaller than \code{:sampled} so Emcee
converges with slightly fewer iterations for the same effective sample size.

The main caveat is that the answer is qualitatively different from
\code{:profile} or \code{:sampled} --- the marginal posterior on $\theta$
will generally be broader because we no longer over-credit a fitted $t$.

\section{Practical considerations}

\subsection{Sampler and threading}

All three model kinds use Emcee with 32 walkers and 500 iterations
($\,32\times500=16\,000$ likelihood evaluations). Emcee is appropriate here
because the posterior is moderately multimodal and the forward model is not
differentiable in a way that supports gradient-based samplers (it has a
\code{try/catch} on \code{DomainError} and mutates buffers, breaking
\code{ForwardDiff}).

Threading: each thread owns its own \code{BasalMixingModel} object, allocated
to length \code{Threads.maxthreadid()} to accommodate Julia's interactive
threads in newer versions.

\subsection{Numerical robustness}

The forward Euler integrator at $\Delta t=0.2\,\mathrm{kyr}$ is marginally
stable; for some proposals the explicit step overshoots and produces
non-positive concentrations, raising \code{DomainError} in
\code{concentration\_to\_age}. The forward function catches these and reports
failure; the \code{@model} retries once at $\Delta t/2$ before giving up and
attaching a vanishingly small likelihood.

\subsection{Outputs and reproducibility}

Each ensemble run writes a snapshot to
\code{results/emcee-chain-<model\_kind>.jld2}, containing the chain, the data
DataFrames, the depth grid, the priors, the sampler and the model kind. The
plotting helpers (\code{plot\_basalmixing\_ensemble.jl}) load this snapshot
and re-derive figures, switching the MAP best-fit re-run on the recorded
\code{model\_kind}.

\section*{References}

\begin{itemize}[leftmargin=*]
  \item Bender, M.~L. et al.\ ``The Contemporary Degassing Rate of $^{40}$Ar
        from the Solid Earth,'' \textit{PNAS} 105 (24), 8232--8237 (2008);
        and the GISP2 measurements used here, archived with the Bender 2010
        data set in this repository.
  \item Goodman, J.\ and Weare, J.\ ``Ensemble Samplers with Affine
        Invariance,'' \textit{Comm.\ App.\ Math.\ Comp.\ Sci.} 5, 65--80 (2010).
  \item Vihola, M.\ ``Robust adaptive Metropolis algorithm with coerced
        acceptance rate,'' \textit{Stat.\ Comput.} 22, 997--1008 (2012).
\end{itemize}

\end{document}
